\section{Examples}

\begin{proposition}[The toolkit of \S18.3: ground states over a finite state space, the last sentence of (18.25)(iii), the counts]
    \label{prop:dse-toolkit}
    \uses{thm:dsb-18-25, def:lo-18-8}
    \lean{MeasureTheory.GibbsMeasure.localGroundEnergy_eq_ciInf, MeasureTheory.GibbsMeasure.localGroundStates_eq_setOf_le, MeasureTheory.GibbsMeasure.exists_pos_localGroundStates_subset, MeasureTheory.GibbsMeasure.localGroundEnergy_eq_of_isOpenPosMeasure, MeasureTheory.GibbsMeasure.cubeSiteSum, MeasureTheory.GibbsMeasure.sqDiagSum, MeasureTheory.GibbsMeasure.cubeBondSum_eq_iff_forall_adj, MeasureTheory.GibbsMeasure.cubeBondSum_onCube_of_spin_refl, MeasureTheory.GibbsMeasure.pairwise_disjoint_faceRestrict_singleton, MeasureTheory.GibbsMeasure.isPhiSymmetric_of_ne, MeasureTheory.GibbsMeasure.CubeTransformation.isSymmetry_of_spin_refl, MeasureTheory.GibbsMeasure.CubeTransformation.IsSymmetry.const_mul, MeasureTheory.GibbsMeasure.exists_gibbsMeasure_family_shift_of_crossingBound_lt_one, MeasureTheory.GibbsMeasure.eventually_exists_gibbsMeasure_family_shift, MeasureTheory.GibbsMeasure.eventually_le_encard_extremePoints, MeasureTheory.GibbsMeasure.eventually_le_encard_extremePoints_G, MeasureTheory.GibbsMeasure.isCubeConfining_univ_of_finite}
    \leanok

    Over a finite state space $G_\varepsilon(\Phi) = \{\Phi_C = m_\Phi\}$ for small $\varepsilon$,
    which turns every ``for $\varepsilon$ small enough'' of \S18.3 into a finite computation; the
    bond and site sums of a cube interaction and their behaviour under the cube isometries; the
    packaged (18.25) exposing the last sentence of (iii), $\theta_{e_k}(\mu_m) = \mu_n$ when
    $r_kG_m = G_n$; the counts $|\operatorname{ex}\mathcal G_\Theta(\beta\Phi)| \geq N$ for an
    $r$-symmetric partition and $|\operatorname{ex}\mathcal G(\beta\Phi)| \geq N$ in general, by
    Corollary (7.29).
\end{proposition}
\begin{proof}
    \leanok
\end{proof}

\begin{example}[Georgii 18.3.1--18.3.2: the Ising ferromagnet and the antiferromagnet in a field]
    \label{ex:dse-ising}
    \uses{prop:dse-toolkit}
    \lean{MeasureTheory.GibbsMeasure.isingFerroCube, MeasureTheory.GibbsMeasure.isingFerroCube_two_values, MeasureTheory.GibbsMeasure.localGroundEnergy_isingFerroCube, MeasureTheory.GibbsMeasure.localGroundStates_isingFerroCube, MeasureTheory.GibbsMeasure.isingFerroG, MeasureTheory.GibbsMeasure.isPStablePartition_isingFerroG, MeasureTheory.GibbsMeasure.isPhiSymmetric_isingFerroG, MeasureTheory.GibbsMeasure.eventually_exists_isingFerro_gibbsMeasure_pair, MeasureTheory.GibbsMeasure.eventually_isingFerro_one_sub_le_measure_zero_mem_latticeOceanPart, MeasureTheory.GibbsMeasure.eventually_two_le_encard_extremePoints_isingFerro, MeasureTheory.GibbsMeasure.isingAntiCube, MeasureTheory.GibbsMeasure.staggerCube, MeasureTheory.GibbsMeasure.neg_four_le_isingAntiCube, MeasureTheory.GibbsMeasure.neg_two_mul_abs_le_isingAntiCube, MeasureTheory.GibbsMeasure.localGroundStates_isingAntiCube, MeasureTheory.GibbsMeasure.eventually_exists_isingAnti_gibbsMeasure_pair}
    \leanok

    18.3.1: $\Phi_C = -\sum_{|i-j|=1}\sigma_i\sigma_j$ takes the values $-4, 0, 4$, so
    $G_\varepsilon(\Phi) = \{(++), (--)\}$ for $0 < \varepsilon < 4$; the partition is stable and
    $\Phi$-symmetric under the spin flip, so for $\beta$ large there are two mutually singular Gibbs
    measures $\mu_\pm \in \mathcal G_\Theta(\beta\Phi)$, invariant under $r_0, r_1, r_2$, exchanged by
    the flip, with $\mu_\pm(\xi^0(G_\pm,\cdot) \neq \emptyset) = 1$ and
    $\mu_\pm(0 \in \xi^0(G_\pm,\cdot)) \geq 1 - 4z(t(G_\varepsilon,\beta\Phi)) \to 1$; hence
    $|\operatorname{ex}\mathcal G_\Theta(\beta\Phi)| \geq 2$, a refinement of (6.9). 18.3.2:
    $\Phi_C = \sum\sigma_i\sigma_j - h\sum_{i\in C}\sigma_i$ with $|h| < 2$ attains $-4$ exactly at the
    two staggered configurations and $-2|h|$ next, so $G_\varepsilon(\Phi)$ is the staggered pair for
    $0 \leq \varepsilon < 4 - 2|h|$; two Gibbs measures invariant under the shifts with $i_1 + i_2$
    even, exchanged by $r_0$, $r_1$, $r_2$, $\theta_{(1,0)}$ and $\theta_{(0,1)}$.
\end{example}

\begin{example}[Georgii 18.3.3: the Gertsik--Dobrushin model]
    \label{ex:dse-gertsik-dobrushin}
    \uses{prop:dse-toolkit}
    \lean{MeasureTheory.GibbsMeasure.gdCube, MeasureTheory.GibbsMeasure.gdCube_apply, MeasureTheory.GibbsMeasure.isSymmetry_of_spin_refl_gdCube, MeasureTheory.GibbsMeasure.isSymmetry_isingSpinFlip_gdCube, MeasureTheory.GibbsMeasure.isSymmetry_sublatticeFlip_gdCube, MeasureTheory.GibbsMeasure.gdCube_onCube_sublatticeFlip, MeasureTheory.GibbsMeasure.gdCube_eq_iff_caseA, MeasureTheory.GibbsMeasure.eventually_two_le_encard_extremePoints_gdCube_caseA, MeasureTheory.GibbsMeasure.gdCube_eq_iff_caseB, MeasureTheory.GibbsMeasure.eventually_exists_gdCube_caseB_gibbsMeasure_pair, MeasureTheory.GibbsMeasure.gdCube_eq_iff_caseC, MeasureTheory.GibbsMeasure.gdGroundC, MeasureTheory.GibbsMeasure.gdFlip, MeasureTheory.GibbsMeasure.isPhiSymmetric_gdGC, MeasureTheory.GibbsMeasure.eventually_four_le_encard_extremePoints_gdCube_caseC, MeasureTheory.GibbsMeasure.gdCube_eq_iff_caseD, MeasureTheory.GibbsMeasure.gdStripe, MeasureTheory.GibbsMeasure.eventually_exists_gdCube_caseD_gibbsMeasure_family, MeasureTheory.GibbsMeasure.eventually_four_le_encard_extremePoints_gdCube_caseD}
    \leanok

    The interaction (18.29) with its value table and symmetries: every spatial cube transformation,
    the spin flip when $h = 0$, and the sublattice flip when $h = J_1 = 0$ --- the sublattice flip
    is not a symmetry in a field, and sends $\Phi^{J_1,J_2,0}$ to $\Phi^{-J_1,J_2,0}$ only, which is
    all Georgii claims. (a) $J_1 > 0$, $J_1 + J_2 > 0$, $h = 0$: two constant ground states,
    $|\operatorname{ex}\mathcal G_\Theta| \geq 2$; (b) $J_1 < 0$, $J_1 < J_2$,
    $|h| < -2J_1 + (J_2 \wedge 0)$: the staggered pair, with the full conclusion of 18.3.2;
    (c) $J_1 = 0$, $J_2 > 0$, $h = 0$: four ground states permuted simply transitively by the
    Klein group of sublattice flips, $|\operatorname{ex}\mathcal G(\beta\Phi)| \geq 4$; (d)
    $|J_1| + J_2 < 0$, $|h| < -J_2 - (J_1 \vee 0)$: the four striped states, four mutually
    singular Gibbs measures, shift invariant in one direction and of period two in the other,
    cycled by $r_0$, $|\operatorname{ex}\mathcal G(\beta\Phi)| \geq 4$. In (c) and (d) the partition
    is not $r$-symmetric, so the count is in $\mathcal G$ rather than $\mathcal G_\Theta$, as in
    Georgii.
\end{example}

\begin{example}[Georgii 18.3.5--18.3.6: the Potts model and the anisotropic Heisenberg ferromagnet]
    \label{ex:dse-potts-heisenberg}
    \uses{prop:dse-toolkit, ex:lo-18-19}
    \lean{MeasureTheory.GibbsMeasure.pottsCube, MeasureTheory.GibbsMeasure.pottsCube_eq_iff, MeasureTheory.GibbsMeasure.pottsShift, MeasureTheory.GibbsMeasure.isPhiSymmetric_pottsG, MeasureTheory.GibbsMeasure.eventually_le_encard_extremePoints_potts, MeasureTheory.GibbsMeasure.eventually_exists_potts_gibbsMeasure_family, MeasureTheory.GibbsMeasure.anisoBond, MeasureTheory.GibbsMeasure.anisoHeisenbergCube, MeasureTheory.GibbsMeasure.anisoBond_eq_one_iff, MeasureTheory.GibbsMeasure.localGroundEnergy_anisoHeisenbergCube, MeasureTheory.GibbsMeasure.anisoG, MeasureTheory.GibbsMeasure.localGroundStates_subset_anisoG, MeasureTheory.GibbsMeasure.anisoRefl, MeasureTheory.GibbsMeasure.eventually_two_le_encard_extremePoints_anisoHeisenberg, MeasureTheory.GibbsMeasure.eventually_exists_anisoHeisenberg_gibbsMeasure_pair}
    \leanok

    18.3.5: on $\mathbb Z/q$ with $\varphi(0) < \varphi(n)$ for $n \neq 0$ (Georgii's
    $\varphi(n) = \varphi(q-n)$ is not needed once the bond term is written $\varphi(x - y)$), the
    local ground states are the $q$ constants, the cyclic shifts of the state space are symmetries
    permuting them, so $|\operatorname{ex}\mathcal G_\Theta(\beta\Phi)| \geq q$ for $\beta$ large.
    18.3.6: the anisotropic bond $(1-\alpha)\langle x,u\rangle\langle y,u\rangle + \alpha\langle x,y\rangle$,
    $0 \leq \alpha < 1$, on the unit sphere with the normalised surface measure has $m_\Phi = -4$ and
    $G_\varepsilon(\Phi)$ inside the two polar caps $\{\pm\langle x,u\rangle \geq 1-\delta\}$ for
    $\varepsilon < (1-\alpha)\delta$, exchanged by the reflection in $u^\perp$, which preserves the
    surface measure: two mutually singular Gibbs measures and
    $|\operatorname{ex}\mathcal G_\Theta(\beta\Phi)| \geq 2$. Not formalised: 18.3.4 (hard-core
    lattice gases, whose interaction takes the value $+\infty$, so the real-valued cube interactions
    of the tree and the positivity used in Step 1 of (18.16) do not apply; Georgii himself notes the
    ad hoc replacement).
\end{example}

\begin{example}[Georgii 18.3.7--18.3.9: the anisotropic antiferromagnet in a field, Shlosman's rotor, double wells]
    \label{ex:dse-rotor-wells}
    \uses{prop:dse-toolkit, ex:dse-potts-heisenberg, ex:dse-ising}
    \lean{MeasureTheory.GibbsMeasure.anisoAntiCube, MeasureTheory.GibbsMeasure.anisoAntiBond, MeasureTheory.GibbsMeasure.anisoAntiCube_eq_cubeBondSum, MeasureTheory.GibbsMeasure.neg_one_le_anisoAntiBond, MeasureTheory.GibbsMeasure.anisoAntiBond_eq_neg_one_iff, MeasureTheory.GibbsMeasure.anisoAntiCube_eq_neg_four_iff, MeasureTheory.GibbsMeasure.localGroundEnergy_anisoAntiCube, MeasureTheory.GibbsMeasure.exists_lt_neg_four_anisoAntiCube, MeasureTheory.GibbsMeasure.anisoAntiG, MeasureTheory.GibbsMeasure.isPStablePartition_anisoAntiG, MeasureTheory.GibbsMeasure.isPhiSymmetric_anisoAntiG, MeasureTheory.GibbsMeasure.eventually_two_le_encard_extremePoints_anisoAnti, MeasureTheory.GibbsMeasure.eventually_exists_anisoAnti_gibbsMeasure_pair, MeasureTheory.GibbsMeasure.exists_pos_localGroundStates_subset_of_isCompact, MeasureTheory.GibbsMeasure.shlosmanCube, MeasureTheory.GibbsMeasure.shlosmanCube_eq_neg_two_iff, MeasureTheory.GibbsMeasure.localGroundEnergy_shlosmanCube, MeasureTheory.GibbsMeasure.shlosmanG, MeasureTheory.GibbsMeasure.isPStablePartition_shlosmanG, MeasureTheory.GibbsMeasure.shlosmanFlip, MeasureTheory.GibbsMeasure.isPhiSymmetric_shlosmanG, MeasureTheory.GibbsMeasure.shlosmanRotation, MeasureTheory.GibbsMeasure.isSymmetry_shlosmanRotation, MeasureTheory.GibbsMeasure.eventually_exists_shlosman_gibbsMeasure_pair, MeasureTheory.GibbsMeasure.eventually_two_le_encard_extremePoints_shlosman, MeasureTheory.GibbsMeasure.doubleWellCube, MeasureTheory.GibbsMeasure.doubleWellCube_eq_min_iff, MeasureTheory.GibbsMeasure.isCompact_sublevel_doubleWellCube, MeasureTheory.GibbsMeasure.doubleWellG, MeasureTheory.GibbsMeasure.isPStablePartition_doubleWellG, MeasureTheory.GibbsMeasure.isPhiSymmetric_doubleWellG, MeasureTheory.GibbsMeasure.isCubeConfining_doubleWellCube, MeasureTheory.GibbsMeasure.eventually_two_le_encard_extremePoints_doubleWell, MeasureTheory.GibbsMeasure.eventually_exists_doubleWell_gibbsMeasure_pair}
    \leanok

    18.3.7, the model (18.34) on the unit sphere with the field in direction $u$: the field
    distributes over the two edges of the square through each site, so the cube interaction is the
    bond sum of $\langle x,u\rangle\langle y,u\rangle + \alpha\langle x_\perp,y_\perp\rangle -
    h(\langle x,u\rangle + \langle y,u\rangle)/2$, whose minimum $-1$ is attained exactly at
    $(u,-u)$ and $(-u,u)$ when $h^2 < 4(1-\alpha^2)$; the two staggered polar-cap patterns are stable,
    exchanged by the reflections, and give $|\operatorname{ex}\mathcal G(\beta\Phi)| \geq 2$ with
    Georgii's invariances. Erratum: Georgii's hypothesis $|h| < 2$ is insufficient for $\alpha > 0$
    --- a tilted (anti)parallel pair has energy $-4\alpha - h^2/(1+\alpha) < -4$ as soon as
    $h^2 > 4(1-\alpha^2)$ (for instance $\alpha = 9/10$, $h = 19/10$), so the staggered pair is then
    not a local ground state; $h^2 < 4(1-\alpha^2)$ is his condition at $\alpha = 0$ and is exactly
    the positive definiteness of the second variation. The transverse-field rotor (18.35), whose
    ground states Georgii computes correctly, is not formalised. 18.3.8, Shlosman's rotor (18.36) on
    the unit circle of an oriented plane, $\cos(\sigma_i - \sigma_j) = \langle\sigma_i,\sigma_j\rangle$: a
    local ground state has orthogonal adjacent spins and equal next-nearest spins, the two
    orientation patterns are stable and exchanged by the spin reversal on the even sublattice, the
    rotations are symmetries preserving each pattern, so two mutually singular Gibbs measures
    invariant under all rotations and the even shifts exist and $|\operatorname{ex}\mathcal
    G(\beta\Phi)| \geq 2$; the variation (18.37) is not formalised. 18.3.9, double wells (18.38) on
    $E = \mathbb R$: the local ground states are the constant configurations at the minimisers of
    $v$, the two wells are stable and exchanged by $\omega \mapsto -\omega$, the confining condition
    (17.18)(iv) in its second alternative holds with $K_\ell = [-\ell-1, \ell+1]$, and
    $|\operatorname{ex}\mathcal G_\Theta(\beta\Phi)| \geq 2$ for any finite symmetric full-support a
    priori measure; Georgii's Lebesgue $\lambda$ enters through his own normalisation
    $\bar\lambda = e^{-4v}\lambda$, and the identification $\gamma^{\beta\bar\Phi,\bar\lambda} =
    \gamma^{\beta\Phi,\lambda}$ is not formalised.
\end{example}

\begin{example}[Georgii 18.3.10: Slawny's $\square$-model, the ground states and the stable partition]
    \label{ex:dse-slawny}
    \uses{prop:dse-toolkit}
    \lean{MeasureTheory.GibbsMeasure.cubeFace, MeasureTheory.GibbsMeasure.faceProd, MeasureTheory.GibbsMeasure.slawnyCube, MeasureTheory.GibbsMeasure.le_slawnyCube, MeasureTheory.GibbsMeasure.slawnyCube_eq_neg_iff, MeasureTheory.GibbsMeasure.slawnyCube_cubeRefl, MeasureTheory.GibbsMeasure.slawnyG, MeasureTheory.GibbsMeasure.slawnyCube_eq_neg_six_subset, MeasureTheory.GibbsMeasure.localGroundEnergy_slawnyCube, MeasureTheory.GibbsMeasure.pairwise_disjoint_faceRestrict_slawnyG, MeasureTheory.GibbsMeasure.slawnyLayerFlip, MeasureTheory.GibbsMeasure.isPhiSymmetric_slawnyG, MeasureTheory.GibbsMeasure.isPStablePartition_slawnyG, MeasureTheory.GibbsMeasure.isCubeConfining_slawnyCube}
    \leanok

    In three dimensions with the plane $P = \{0,1\}$: the interaction (18.39) is minimal iff each
    face of the cube carries an even number of $+1$ spins; if two vertically adjacent spins agree
    at one site of a ground state they agree everywhere, so the ground states split into the two
    patterns $G_1, G_2$, which are $P$-stable although their projections on the top and bottom
    faces coincide, and are exchanged by the layer flip $\omega \mapsto ((-1)^{i_3}\omega_i)$; the
    interaction is bounded, hence confining. The conclusion of Georgii and (18.40)--(18.41) are not
    formalised: they need Theorem (18.25) and \S18.1 in dimension $d \geq 3$ with the plane $P$,
    whereas the tree's ocean part and (18.25) are two-dimensional; the model-level input is complete.
\end{example}
